\section{Appendix}
\label{app:extra}

\subsection{Mechanistic Formalization}
\label{app:theory-bridge}
This appendix gives the assumptions and proofs for the three compact propositions stated in Setup.
The goal is not a general theorem about deep networks under shift.
It is to formalize the selector mismatch created by suppressive weighting, explain why the validation-loss guardrail helps, and give a principled bridge from the empirical time courses to the transient-versus-persistent split.

\paragraph{Proposition 1 (Local lure under suppressive weighting).}
Assume a suppressive objective whose high-loss slope ratio is $\alpha \in (0,1)$.
Let $A_c>0$ denote core improvement mass and $A_t>0$ denote tail damage mass over a short interval.
For local loss tradeoffs in which the core improvement is on examples below the clipping threshold and the tail damage is on examples above it, there is a non-empty region in which the suppressive proxy improves while standard risk worsens, namely $\alpha A_t < A_c < A_t$.

\emph{Proof.}
For soft clipping with threshold $\tau$ and slope $\alpha \in (0,1]$,
\[
\tilde{\ell} = \ell - (1-\alpha)(\ell-\tau)_+,
\]
so clipping removes upper-tail loss mass directly at the per-example level.
When $\tau=\mathrm{VaR}_q(\ell)$, the same monotonicity gives a compressed proxy tail
\[
\mathrm{CVaR}_q(\tilde{\ell}) = \tau + \alpha\big(\mathrm{CVaR}_q(\ell)-\tau\big),
\]
which is why the proxy can look improved even when the underlying standard-loss tail has not.

\paragraph{Local lure condition.}
Let $A_c>0$ denote core improvement mass and $A_t>0$ denote tail damage mass over a short interval.
The short-interval approximation assumes that the core improvement mass lies below the clipping threshold, where the proxy slope is $1$, and the tail damage mass lies above it, where the proxy slope is $\alpha$.
Then standard risk and suppressive proxy move as
\[
\Delta R \approx -A_c + A_t,
\qquad
\Delta P \approx -A_c + \alpha A_t.
\]
Whenever
\[
\alpha A_t < A_c < A_t,
\]
the proxy still improves while the standard risk worsens.
So the local lure region is non-empty whenever $\alpha<1$.
It widens as $\alpha$ decreases, which is why stronger suppression creates more room for proxy-only selection to move toward less reliable checkpoints.
\hfill$\square$

\paragraph{Proposition 2 (Baseline-referenced guardrail bound).}
Assume the guardrail enforces
\[
L_{\mathrm{val}}(\hat\theta_{\mathrm{guard}}) \le \rho L_{\mathrm{val}}(\hat\theta_{\mathrm{base}})
\]
and that validation and deployment losses satisfy
\[
\sup_t \left|L_{\mathrm{test}}(\theta_t)-L_{\mathrm{val}}(\theta_t)\right| \le \varepsilon
\]
over the candidate suppressive checkpoints and also for the baseline checkpoint $\hat\theta_{\mathrm{base}}$.
Then
\[
L_{\mathrm{test}}(\hat\theta_{\mathrm{guard}})
\le
\rho\,L_{\mathrm{test}}(\hat\theta_{\mathrm{base}})
+
(\rho+1)\varepsilon.
\]

\emph{Proof.}
By the uniform closeness assumption,
\[
L_{\mathrm{test}}(\hat\theta_{\mathrm{guard}})
\le
L_{\mathrm{val}}(\hat\theta_{\mathrm{guard}})+\varepsilon
\le
\rho\,L_{\mathrm{val}}(\hat\theta_{\mathrm{base}})+\varepsilon.
\]
Applying the same bound to the baseline checkpoint gives
\[
L_{\mathrm{val}}(\hat\theta_{\mathrm{base}})
\le
L_{\mathrm{test}}(\hat\theta_{\mathrm{base}})+\varepsilon,
\]
so
\[
L_{\mathrm{test}}(\hat\theta_{\mathrm{guard}})
\le
\rho\,L_{\mathrm{test}}(\hat\theta_{\mathrm{base}})
+
(\rho+1)\varepsilon.
\]
This does not make the guardrail optimal, but it does show why the rule is not arbitrary: the budget caps deployment loss relative to baseline up to validation-to-test mismatch, while proxy-only selection has no analogous guarantee without assuming that the suppressive proxy stays aligned with standard reliability metrics.
\hfill$\square$

\paragraph{Proposition 3 (Transient versus persistent in a two-phase toy).}
Let $H_A\ge 0$ denote the excess standard-loss harm accumulated during an early suppressive phase, and let $H_B\ge 0$ denote the later recovery.
Then the net hazard is persistent when $H_A > H_B$ and transient when $H_B \ge H_A$.

\emph{Proof.}
By construction, the two-phase toy decomposes the total standard-loss effect into an early excess-harm term and a later recovery term.
The sign of the net effect is therefore the sign of $H_A-H_B$.
If $H_A > H_B$, the run remains worse than baseline at the matched horizon; if $H_B \ge H_A$, the excess harm has been offset and the matched-horizon hazard disappears.
\hfill$\square$

Figure~\ref{fig:theory-bridge-app} illustrates both the local lure region and this two-phase transient-versus-persistent picture.

\begin{figure}[h]
  \centering
  \includegraphics[width=\columnwidth]{fig_theory_bridge}
  \caption{
    Mechanistic formalization.
    Left: the local lure region where the suppressive proxy still improves while standard risk worsens.
    Right: a two-phase toy showing how later recovery yields a transient hazard while entrenched suppression yields a persistent one.
  }
  \label{fig:theory-bridge-app}
\end{figure}

\subsection{Persistence Bridge via $R_w(t)$}
Table~\ref{tab:rw-persistence-bridge} and Figure~\ref{fig:rw-persistence-bridge-app} provide the empirical bridge for the transient-versus-persistent taxonomy.
The main discriminator is not a single endpoint value of $R_w$, but whether suppressive orientation relaxes or becomes more entrenched over training, so we use the early-to-late trend in $R_w$ together with the tail-core clip-gap trend as the primary persistence surrogates.

\begin{table}[h]
  \centering
  \small
  \resizebox{\columnwidth}{!}{\begin{tabular}{lccccc}
  \toprule
  Suite & Early $R_w$ & Late $R_w$ & $\Delta R_w$ & $\Delta$clip gap & Trend \\
  \midrule
  Camelyon17 & 0.277 $\pm$ 0.003 & 0.187 $\pm$ 0.002 & -0.091 $\pm$ 0.003 & 0.095 $\pm$ 0.004 & entrenches \\
  CelebA & 0.196 $\pm$ 0.013 & 0.298 $\pm$ 0.011 & 0.102 $\pm$ 0.017 & -0.087 $\pm$ 0.013 & relaxes \\
  \bottomrule
\end{tabular}
}
  \caption{
    Time-course bridge for transient versus persistent hazards.
    CelebA relaxes over training ($\Delta R_w > 0$ and decreasing tail-core clip gap), while Camelyon17 becomes more entrenched ($\Delta R_w < 0$ and increasing clip gap).
  }
  \label{tab:rw-persistence-bridge}
\end{table}

\begin{figure}[h]
  \centering
  \includegraphics[width=\columnwidth]{fig_rw_persistence_bridge}
  \caption{
    Early-versus-late suppressive orientation on the two main time-course anchors.
    The trajectories make the transient-versus-persistent split visually explicit:
    CelebA relaxes, while Camelyon17 entrenches.
  }
  \label{fig:rw-persistence-bridge-app}
\end{figure}

\subsection{Core-Case Uncertainty}
Table~\ref{tab:core-case-uncertainty} gives paired $95\%$ bootstrap intervals for the three seed-matched hazard anchors used most heavily in the main text.
For the Camelyon17 suites, the supporting diagnostic column is teacher-difficulty tail CVaR; for CivilComments, the reliability column is official WILDS worst-group loss.

\begin{table}[h]
  \centering
  \small
  \resizebox{\columnwidth}{!}{\begin{tabular}{lccc}
  \toprule
  Case & $\Delta$Held-out loss & $\Delta$Tail / reliability & $\Delta$ECE \\
  \midrule
  Camelyon17 ERM & +0.055 [+0.045, +0.066] & +5.50 [+5.00, +6.04] & +0.032 [+0.030, +0.033] \\
  Camelyon17 Finetune & +0.619 [+0.331, +0.926] & +6.23 [+3.49, +9.07] & +0.021 [+0.018, +0.025] \\
  CivilComments & +0.104 [+0.100, +0.107] & +1.829 [+1.785, +1.871] & +0.070 [+0.068, +0.071] \\
  \bottomrule
\end{tabular}
}
  \caption{
    Paired bootstrap intervals for the main hazard anchors.
    These intervals complement the selected-checkpoint deltas in Table~\ref{tab:core-cases} by making the uncertainty on held-out loss, reliability, and calibration explicit.
  }
  \label{tab:core-case-uncertainty}
\end{table}

\subsection{Tail-Bank Protocol and Per-Bank Robustness}
\label{app:tail-bank-robustness}
Because the teacher-tail stress diagnostic is supporting fixed-bank evidence on the vision anchors, we make the bank construction explicit here.
The two fixed teacher-difficulty banks are produced once per stack by \texttt{src/scripts/build\_conf\_eval\_banks.py}.
For Camelyon17, bank A and bank B come from two frozen ERM teacher runs (seed 0 and seed 1, recorded in \texttt{artifacts/partitions\_eval/camelyon17\_resnet50/teacher\_difficulty/meta.json}).
Within each bank, examples are assigned to $K=64$ quantile cells by per-example teacher cross-entropy loss, with a minimum evaluation cell size of 20 enforced downstream by \texttt{src/scripts/phase1\_eval.py}.
Those bank assignments are fixed before selector comparison and are not redefined per checkpoint.
Table~\ref{tab:camelyon-per-bank-robustness-app} and Table~\ref{tab:camelyon-finetune-per-bank-robustness-app} then rerun the selected-checkpoint comparison separately on banks A and B using \texttt{src/scripts/analyze\_per\_bank\_robustness.py}.
For Camelyon17 ERM, the harmful direction holds cleanly on both banks at P95, P97, and P99.
For finetuning, the selected-checkpoint bankwise direction remains at P95, P97, and P99 as well: the harmful mean tail direction holds on both banks, with slightly weaker decoupling fractions at P95 on both banks and P97 on bank B.

\begin{table}[h]
  \centering
  \small
  \resizebox{\columnwidth}{!}{\begin{tabular}{llccccc}
  \toprule
  Dataset & Bank & Regime & $\Delta$Proxy & $\Delta$Tail & Proxy improve frac & Decoupling frac \\
  \midrule
  Camelyon17 & A & P95 & -0.297 & +5.501 & 1.00 & 1.00 \\
  Camelyon17 & A & P97 & -0.267 & +5.155 & 1.00 & 1.00 \\
  Camelyon17 & A & P99 & -0.151 & +1.765 & 1.00 & 1.00 \\
  Camelyon17 & B & P95 & -0.258 & +5.501 & 1.00 & 1.00 \\
  Camelyon17 & B & P97 & -0.170 & +5.155 & 1.00 & 1.00 \\
  Camelyon17 & B & P99 & -0.029 & +1.765 & 1.00 & 1.00 \\
  \bottomrule
\end{tabular}
}
  \caption{
    Camelyon17 ERM per-bank robustness for the teacher-difficulty tail metric.
    Rows report mean proxy and tail deltas relative to the ERM baseline separately on banks A and B.
    The harmful teacher-difficulty tail direction holds on both banks, so the main P95 result is not driven by averaging across a single favorable bank.
    For ERM, the bankwise tail deltas coincide after rounding in the saved artifact, while proxy deltas differ by bank.
  }
  \label{tab:camelyon-per-bank-robustness-app}
\end{table}

\begin{table}[h]
  \centering
  \small
  \resizebox{\columnwidth}{!}{\begin{tabular}{llccccc}
  \toprule
  Dataset & Bank & Regime & $\Delta$Proxy & $\Delta$Tail & Proxy improve frac & Decoupling frac \\
  \midrule
  Camelyon17 finetune selected & A & P95 & -0.605 & +6.430 & 1.00 & 0.90 \\
  Camelyon17 finetune selected & A & P97 & -0.528 & +5.409 & 1.00 & 1.00 \\
  Camelyon17 finetune selected & A & P99 & -0.359 & +5.744 & 1.00 & 1.00 \\
  Camelyon17 finetune selected & B & P95 & -0.458 & +6.022 & 1.00 & 0.90 \\
  Camelyon17 finetune selected & B & P97 & -0.391 & +5.484 & 1.00 & 0.90 \\
  Camelyon17 finetune selected & B & P99 & -0.261 & +5.653 & 1.00 & 1.00 \\
  \bottomrule
\end{tabular}
}
  \caption{
    Camelyon17 finetune selected-checkpoint per-bank robustness for the teacher-difficulty tail metric.
    The harmful tail direction remains on both banks at P95, P97, and P99.
    P95 and P97 are slightly less uniform than the ERM anchor on one of the two bankwise consistency fractions, while P99 remains fully harmful on both banks.
  }
  \label{tab:camelyon-finetune-per-bank-robustness-app}
\end{table}

\subsection{Persistent-Case Fixed-Horizon Support}
\label{app:persistence-support}
Table~\ref{tab:persistence-support-app} exposes the matched fixed-horizon comparison for every case we still label persistent in the main paper.
For the classification anchors, the matched horizons are epoch 30 (Camelyon17 ERM and frozen CivilComments) and epoch 10 (Camelyon17 finetuning), matching the main-text setup.
For ACSIncome, all runs share the same 50-epoch training horizon, so the fixed-horizon comparison uses the final epoch.
The goal is not to claim identical magnitudes between selected and fixed views; it is to show that the harmful loss / reliability direction survives the matched-horizon check for every row that retains the persistent label.
For Camelyon17 finetuning, the fixed10 teacher-difficulty tail diagnostic is markedly more outlier-sensitive than the selected-checkpoint view, so we use it only as a qualitative persistence check rather than as a clean magnitude claim.

\begin{table}[h]
  \centering
  \small
  \resizebox{\columnwidth}{!}{\begin{tabular}{llcrrrrl}
  \toprule
  Case & Fixed horizon & $n$ & $\Delta$Loss (sel.) & $\Delta$Loss (fixed) & $\Delta$Rel. (sel.) & $\Delta$Rel. (fixed) & Verdict \\
  \midrule
  Camelyon17 ERM & ep30 & 10 & +0.055 & +0.053 & +5.501 & +5.434 & Persistent \\
  Camelyon17 Finetune & ep10 & 10 & +0.619 & +0.491 & +6.226 & mean $>0^{\dagger}$ & Persistent \\
  Frozen CivilComments & ep30 & 10 & +0.104 & +0.124 & +1.829 & +1.923 & Persistent \\
  ACSIncome & ep50 & 10 & +0.027 & +0.017 & +0.376 & +0.284 & Persistent \\
  \bottomrule
\end{tabular}
}
  \caption{
    Fixed-horizon support for every row still labeled persistent.
    ``Selected'' is the main-text proxy-selected comparison; ``fixed'' evaluates the same baseline/suppressive pair at the matched horizon instead of at the selected epoch.
    For ACSIncome, the matched fixed horizon is the common final epoch 50.
    $\dagger$ The finetune fixed10 teacher-difficulty tail diagnostic is outlier-sensitive: a few seeds explode at epoch 10, so we use that entry only as a qualitative persistence check rather than as a clean magnitude comparison.
  }
  \label{tab:persistence-support-app}
\end{table}

\subsection{Suppression-Strength Dose Response on Camelyon}
\label{app:camelyon-dose-response}
Table~\ref{tab:camelyon-dose-response-app} adds a direct suppression-strength stress test on the two strongest Camelyon anchors.
For ERM, weakening the clipping threshold from P95 to P97 to P99 reduces the baseline-referenced selected-checkpoint harm smoothly in both held-out loss and teacher-difficulty tail CVaR, so the main P95 result is not a knife-edge artifact of one particular threshold.
For finetuning, the selected-checkpoint harm remains large across all three regimes, and the bankwise selected-checkpoint comparison continues to show harmful tail movement on both banks (Table~\ref{tab:camelyon-finetune-per-bank-robustness-app}).
The matched fixed10 finetune view is materially noisier, so we do not present the finetune dose-response as a monotone law; the supportable conclusion is narrower: once suppressive finetuning is allowed to select by its own proxy, the selected-checkpoint hazard persists beyond a single P95 setting.

\begin{table}[h]
  \centering
  \small
  \resizebox{\columnwidth}{!}{\begin{tabular}{llrrr}
\toprule
Suite & Regime & Frac clip & $\Delta$Loss$\downarrow$ & $\Delta$Tail$\downarrow$ \\
\midrule
Camelyon17 ERM & P95 & 0.089 & +0.055 & +5.50 \\
 & P97 & 0.076 & +0.052 & +5.16 \\
 & P99 & 0.037 & +0.008 & +1.76 \\
Camelyon17 finetune & P95 & 0.068 & +0.619 & +6.23 \\
 & P97 & 0.060 & +0.772 & +5.45 \\
 & P99 & 0.043 & +0.551 & +5.70 \\
\bottomrule
\end{tabular}
}
  \caption{
    Camelyon suppression-strength dose response on the selected checkpoints.
    Frac clip reports the mean validation clipping fraction at the selected epoch.
    In the baseline-referenced selected-checkpoint workflow view, ERM loss / tail harm weakens as clipping relaxes from P95 to P97 to P99.
    Table~\ref{tab:camelyon-erm-selector-dose} separately reports the same-trajectory selector-dose diagnostic.
    On finetuning, the selected-checkpoint hazard remains large across all three tested clipping thresholds, so P95 should be read as the designated strong stress-test setting rather than as a unique knife-edge.
  }
  \label{tab:camelyon-dose-response-app}
\end{table}

\subsection{Guardrail Uncertainty}
\label{app:guardrail}
Table~\ref{tab:guardrail-uncertainty} reports paired bootstrap intervals for the improvement achieved by the guardrail relative to proxy-only selection.
The intervals are computed on seed-matched proxy-only versus guardrail selections.

\begin{table}[h]
  \centering
  \small
  \resizebox{\columnwidth}{!}{\begin{tabular}{llccc}
  \toprule
  Suite & Budget & Loss reduction & Tail reduction & Acc. change \\
  \midrule
  Camelyon17 ERM & 1.10x & +0.055 [+0.046,+0.066] & +5.501 [+5.001,+6.022] & -0.001 [-0.002,+0.001] \\
  Camelyon17 ERM & 1.25x & +0.029 [+0.016,+0.044] & +2.988 [+1.873,+4.209] & -0.003 [-0.004,-0.001] \\
  Camelyon17 ERM & 1.50x & +0.003 [+0.000,+0.010] & +0.177 [+0.000,+0.532] & +0.000 [+0.000,+0.001] \\
  Camelyon17 Finetune & 1.10x & +0.626 [+0.310,+0.956] & +6.382 [+3.712,+9.171] & -0.024 [-0.074,+0.027] \\
  Camelyon17 Finetune & 1.25x & +0.551 [+0.260,+0.873] & +5.767 [+2.684,+8.879] & -0.015 [-0.057,+0.027] \\
  Camelyon17 Finetune & 1.50x & +0.301 [+0.061,+0.580] & +4.234 [+1.036,+7.813] & -0.025 [-0.058,+0.006] \\
  \bottomrule
\end{tabular}
}
  \caption{
    Paired bootstrap intervals for proxy-only $\rightarrow$ guardrail changes.
    Positive loss/tail reductions mean the guardrail removes harm relative to proxy-only selection.
  }
  \label{tab:guardrail-uncertainty}
\end{table}

\subsection{Guardrail Budget Sweep Detail}
Table~\ref{tab:guardrail-budget-app} retains the full budget sweep that sits behind the main-text selector comparison.
It shows why the main paper centers the selector comparison on the $1.25\times$ operating point as a middle diagnostic budget rather than presenting it as a universal constant.

\begin{table}[h]
  \centering
  \small
  \resizebox{\columnwidth}{!}{\begin{tabular}{llcccc}
  \toprule
  Suite & Rule & Accept rate & $\Delta$Held-out loss$\downarrow$ & $\Delta$Held-out acc$\uparrow$ & $\Delta$Tail CVaR$\downarrow$ \\
  \midrule
  Camelyon17 ERM & Proxy-only & 1.00 & +0.055 & +0.001 & +5.501 \\
   & 1.10x budget & 0.00 & +0.000 & +0.000 & +0.000 \\
   & 1.25x budget & 0.70 & +0.026 & -0.002 & +2.513 \\
   & 1.50x budget & 1.00 & +0.052 & +0.001 & +5.324 \\
  Camelyon17 Finetune & Proxy-only & 1.00 & +0.619 & +0.025 & +6.226 \\
   & 1.10x budget & 0.10 & +0.000 & +0.000 & -0.156 \\
   & 1.25x budget & 0.40 & +0.068 & +0.010 & +0.460 \\
   & 1.50x budget & 0.70 & +0.318 & +0.000 & +1.992 \\
  \bottomrule
\end{tabular}
}
  \caption{
    Full budget sweep on the two Camelyon17 anchor suites.
    The $1.10\times$ budget is essentially a veto, $1.50\times$ is materially looser, and $1.25\times$ is the middle illustrative diagnostic budget used in the main text.
  }
  \label{tab:guardrail-budget-app}
\end{table}

\subsection{Adaptive Baseline-Only Guardrail Sensitivity}
\label{app:adaptive-guardrail}
To address reviewer-facing concerns about $\rho$, we also checked a fully validation-only adaptive heuristic.
Instead of using the fixed $1.25\times$ budget from Equation~\ref{eq:guardrail}, it sets the admissible threshold to $\mu+2\sigma$ from the final stable baseline-loss window and then applies proxy-based selection only within that feasible set.
Table~\ref{tab:adaptive-guardrail-app} shows that the result is conservative rather than contradictory: on Camelyon17 ERM it becomes a pure veto, while on finetuning it narrows but does not eliminate proxy-only damage.
We therefore treat adaptive baseline-only budgets as a sensitivity analysis, not as a replacement for the illustrative fixed-budget operating point used in the main selector comparison.

\begin{table}[h]
  \centering
  \small
  \resizebox{\columnwidth}{!}{\begin{tabular}{lccccc}
  \toprule
  Suite & Selector & Admit rate & $\Delta$Held-out loss$\downarrow$ & $\Delta$Held-out acc$\uparrow$ & $\Delta$Tail$\downarrow$ \\
  \midrule
  Camelyon17 ERM & Baseline & 1.00 & +0.000 & +0.000 & +0.000 \\
   & Proxy-only & 1.00 & +0.055 & +0.001 & +5.501 \\
   & Adaptive $\mu+2\sigma$ & 0.00 & +0.000 & +0.000 & +0.000 \\
  Camelyon17 Finetune & Baseline & 1.00 & +0.000 & +0.000 & +0.000 \\
   & Proxy-only & 1.00 & +0.619 & +0.025 & +6.226 \\
   & Adaptive $\mu+2\sigma$ & 0.90 & +0.482 & +0.031 & +5.190 \\
  \bottomrule
\end{tabular}
}
  \caption{
    Baseline-only adaptive guardrail sensitivity on the Camelyon17 selector anchors.
    The adaptive threshold is set from final baseline validation-loss fluctuations alone ($\mu+2\sigma$; last five epochs on 30-epoch ERM and last three epochs on 10-epoch finetuning).
    Deltas are measured relative to the baseline-selected checkpoint.
    The heuristic is conservative on ERM and still directionally safer than proxy-only on finetuning, but it is not as effective as the illustrative fixed $1.25\times$ operating point from the main text.
  }
  \label{tab:adaptive-guardrail-app}
\end{table}

\subsection{Public-Practice Audit of Scalarized or Under-Specified Selection}
\label{app:tooling-audit}
We audited $10$ adjacent published papers with released artifacts, spanning benchmark/selection artifacts and shaped-objective or robust-learning repos (DomainBed, WILDS, SWAD, SubpopBench, ERM++, group DRO \citep{sagawa2020groupdro}, To Smooth or Not? \citep{wei2022smoothnoisy}, LS-PLL \citep{gong2024lspll}, Symmetric Cross Entropy \citep{wang2019symmetriccrossentropy}, and SGN \citep{englesson2024robustclassificationregression}).
The sample is illustrative only, chosen to span adjacent released benchmark/selection artifacts and suppressive-or-robust objective repos rather than to estimate frequency or representativeness.
Table~\ref{tab:tooling-audit} shows the pattern: $8/10$ expose checkpoint choice through a single best-model scalar, and none exposes a default reliability-aware multi-metric selector or checks alignment with held-out loss, calibration, or group-sensitive reliability.

\begin{table}[h]
  \centering
  \small
  \resizebox{\columnwidth}{!}{\begin{tabular}{@{}p{1.75cm}p{1.25cm}p{3.55cm}cc@{}}
  \toprule
  Artifact & Family & Exposed best-model rule & \shortstack[c]{Scalarized\\default?} & \shortstack[c]{Loss/calib\\cross-check?} \\
  \midrule
  DomainBed & DG & top \texttt{val\_acc} & Yes & No \\
  \shortstack[l]{WILDS\\ runner} & shift & \shortstack[l]{best epoch by\\\texttt{val\_metric}} & Yes & No \\
  group\_DRO & \shortstack[l]{robust\\rwgt.} & \shortstack[l]{worst-group /\\avg.\ val acc} & Yes & No \\
  SWAD & DG & in-domain val acc & Yes & No \\
  SubpopBench & subgroup & \shortstack[l]{val worst-group /\\worst-class acc} & Yes & No \\
  ERM++ & DG & highest val acc & Yes & No \\
  \shortstack[l]{To Smooth or\\ Not?} & \shortstack[l]{label\\smooth.} & held-out acc & Yes & No \\
  LS-PLL & \shortstack[l]{partial\\labels} & test acc & Yes & No \\
  \shortstack[l]{Symmetric\\ CE} & \shortstack[l]{noisy\\labels} & \shortstack[l]{epochwise\\\texttt{val\_loss}} & Yes & No \\
  SGN & \shortstack[l]{noisy\\labels} & \shortstack[l]{no default\\best-model rule} & No & No \\
  \bottomrule
\end{tabular}
}
  \caption{
    Compact audit of checkpoint-selection defaults in adjacent published papers with released artifacts.
    The sample spans benchmark/selection artifacts and shaped-objective or robust-learning repos.
    This is a purposive relevance audit rather than a prevalence estimate; the $8/10$ count should not be read as a population rate.
    The common pattern is not a specific metric, but that checkpoint choice is scalarized or left under-specified by default rather than exposed through a reliability-aware multi-metric selector.
  }
  \label{tab:tooling-audit}
\end{table}

\subsection{End-to-End CivilComments Support}
\label{app:civilcomments-text-e2e}
As a targeted scope check, we also reran CivilComments end-to-end with DistilBERT finetuning on an official-group-aware skewed validation surface emphasizing minority-identity positive examples.
Table~\ref{tab:civilcomments-text-e2e} reports the ten-seed selector comparison.
Relative to the baseline-selected checkpoint, proxy-only soft clipping worsens overall loss by $+0.155$, official WILDS worst-group loss by $+1.007$, and ECE by $+0.0533$; validation-loss-only remains worse than baseline, while the $1.25\times$ guardrail acts as a pure veto on all $10/10$ seeds.

Unlike the frozen-embedding CivilComments result, fixed-horizon official worst-group persistence is mixed in this end-to-end text stack: at epoch 10, soft clipping is still worse on overall loss and ECE, but only $5/10$ seeds are worse on official worst-group loss and worst-group accuracy is often higher ($8/10$).
So this extension stays in the appendix; its value is narrower, namely that the selected-checkpoint loss / calibration hazard survives end-to-end text finetuning even when the fixed-horizon worst-group signal is mixed, consistent with early suppressive $R_w$ relaxing toward $1$ later in training (Figure~\ref{fig:civilcomments-text-mechanism-app}).

\begin{table}[h]
  \centering
  \small
  \resizebox{0.92\columnwidth}{!}{\begin{tabular}{lccccc}
\toprule
Selector & $\Delta$Acc & $\Delta$Loss & $\Delta$WG Acc & $\Delta$WG Loss & $\Delta$ECE \\
\midrule
Proxy-only & -0.0046 & +0.155 & +0.0330 & +1.007 & +0.0533 \\
Val-loss-only & -0.0022 & +0.096 & -0.0583 & +0.919 & +0.0471 \\
Guardrail $1.25\times$ & +0.0000 & +0.000 & +0.0000 & +0.000 & +0.0000 \\
\bottomrule
\end{tabular}
}
  \caption{
    End-to-end CivilComments selector comparison on the official-group-aware validation surface.
    Deltas are measured relative to the baseline-selected checkpoint.
    Proxy-only soft clipping worsens overall loss and official WILDS worst-group loss, while the $1.25\times$ guardrail acts as a pure veto and falls back to baseline on all ten seeds.
    Rows are backed by the saved selector-summary artifact listed in the supplement ledger.
  }
  \label{tab:civilcomments-text-e2e}
\end{table}

\begin{figure}[h]
  \centering
  \includegraphics[width=\columnwidth]{civilcomments_text_mechanism_wgminor_10s_20260331_ext}
  \caption{
    End-to-end CivilComments mechanism on the official-group-aware validation surface.
    Left: early $R_w$ is suppressive before relaxing toward $1$ later in training.
    Right: the tail-core clip-exposure gap remains large, showing that the selected checkpoints still sit in a genuinely suppressive regime.
  }
  \label{fig:civilcomments-text-mechanism-app}
\end{figure}

\subsection{ACSIncome Huber Dose-Response}
\label{app:acsincome-huber}
On the same ACSIncome regression stack, Huber loss at the P95 residual threshold provides a milder suppressive objective ($R_w = 0.53$, vs.\ $0.41$ for winsorization).
The selection hazard is weaker and shifts toward the tail channel: proxy-only selection worsens raw held-out MSE on $3/10$ seeds and tail MSE on $9/10$ seeds.
The mean $\Delta$\,tail MSE is $+0.051$ with $95\%$ CI $[+0.028, +0.073]$ (excludes zero), while the mean $\Delta$\,held-out MSE is slightly negative at $-0.0031$ with $95\%$ CI $[-0.0053, -0.0008]$.
This is consistent with Proposition~1's dose-response picture: Huber preserves a weaker but non-zero tail signal, and the ordering $R_w^{\text{wins}} = 0.41 < R_w^{\text{huber}} = 0.53$ maps to a much stronger winsorization hazard.

\subsection{Seed-Matched Lure Checks}
The main paper reports the strongest seed-level counts in prose.
Table~\ref{tab:seed-lure-dominance} gives the full dominance counts, while Table~\ref{tab:seed-lure-paired} gives paired mean deltas with bootstrap intervals.
For Camelyon17, the reliability diagnostic is teacher-difficulty tail CVaR; for CivilComments, it is official WILDS worst-group loss.

\begin{table}[h]
  \centering
  \small
  \resizebox{\columnwidth}{!}{\begin{tabular}{lccccc}
  \toprule
  Suite & Proxy better & Proxy$\land$val loss worse & Proxy$\land$test loss worse & Proxy$\land$reliability worse & Proxy$\land$ECE worse \\
  \midrule
  Camelyon17 ERM & \textbf{10/10} & \textbf{10/10} & \textbf{10/10} & \textbf{10/10} & \textbf{10/10} \\
  Camelyon17 Finetune & 8/10 & 8/10 & 7/10 & 7/10 & 8/10 \\
  CivilComments & \textbf{10/10} & \textbf{10/10} & \textbf{10/10} & \textbf{10/10} & \textbf{10/10} \\
  \bottomrule
\end{tabular}
}
  \caption{
    Seed-level lure counts.
    ``Reliability worse'' means higher teacher-difficulty tail CVaR on Camelyon17 and higher official WILDS worst-group loss on CivilComments.
  }
  \label{tab:seed-lure-dominance}
\end{table}

\begin{table}[h]
  \centering
  \small
  \resizebox{\columnwidth}{!}{\begin{tabular}{llccc}
  \toprule
  Suite & Metric & Mean & 95\% CI low & 95\% CI high \\
  \midrule
  Camelyon17 ERM & proxy clip & -0.0709 & -0.0769 & -0.0659 \\
  Camelyon17 ERM & validation loss & +0.0891 & +0.0762 & +0.1034 \\
  Camelyon17 ERM & held-out loss & +0.0553 & +0.0455 & +0.0658 \\
  Camelyon17 ERM & tail CVaR & +5.5015 & +5.0025 & +6.0403 \\
  Camelyon17 ERM & ECE & +0.0321 & +0.0304 & +0.0334 \\
  Camelyon17 ERM & Brier & +0.0027 & +0.0014 & +0.0038 \\
  Camelyon17 ERM & held-out acc & +0.0010 & -0.0007 & +0.0025 \\
  Camelyon17 Finetune & proxy clip & -0.0141 & -0.0377 & +0.0119 \\
  Camelyon17 Finetune & validation loss & +0.1229 & +0.0754 & +0.1733 \\
  Camelyon17 Finetune & held-out loss & +0.6185 & +0.3311 & +0.9258 \\
  Camelyon17 Finetune & tail CVaR & +6.2264 & +3.4939 & +9.0689 \\
  Camelyon17 Finetune & ECE & +0.0214 & +0.0178 & +0.0249 \\
  Camelyon17 Finetune & Brier & +0.0006 & -0.0033 & +0.0051 \\
  Camelyon17 Finetune & held-out acc & +0.0251 & -0.0218 & +0.0700 \\
  CivilComments & proxy clip & -0.0143 & -0.0151 & -0.0135 \\
  CivilComments & validation loss & +0.0609 & +0.0583 & +0.0638 \\
  CivilComments & held-out loss & +0.1037 & +0.1002 & +0.1071 \\
  CivilComments & ECE & +0.0701 & +0.0685 & +0.0714 \\
  CivilComments & Brier & +0.0138 & +0.0134 & +0.0143 \\
  CivilComments & held-out acc & -0.0040 & -0.0045 & -0.0034 \\
  CivilComments & WILDS wg loss & +1.8292 & +1.7853 & +1.8710 \\
  CivilComments & WILDS wg acc & -0.1474 & -0.1545 & -0.1394 \\
  CivilComments & WILDS wg Brier & +0.3003 & +0.2945 & +0.3061 \\
  CivilComments & WILDS wg ECE & +0.3041 & +0.2952 & +0.3126 \\
  \bottomrule
\end{tabular}
}
  \caption{
    Paired mean deltas for the main lure checks.
    ``Reliability'' denotes teacher-difficulty tail CVaR on Camelyon17 and official WILDS worst-group loss on CivilComments; calibration is reported through both ECE and Brier.
  }
  \label{tab:seed-lure-paired}
\end{table}

\subsection{Camelyon Standard-Metric, Selector-Dose, and Leave-One-Hospital Diagnostics}
\label{app:camelyon-selector-diagnostics}
Table~\ref{tab:camelyon-selector-diagnostics} adds two targeted diagnostics for the Camelyon anchors.
The workflow rows show that the Camelyon reliability signal is not carried only by teacher-difficulty tail CVaR: on ERM, held-out loss, ECE, and Brier all worsen; on finetuning, held-out loss and validation-side ECE worsen while Brier is neutral.
The same-trajectory rows compare proxy-best with validation-loss-best inside the same suppressive trajectory, with no fallback.
This selector-only contrast is clean on Camelyon17 ERM and mixed on finetuning, so the finetuning claim remains primarily baseline-referenced workflow evidence.

\begin{table}[h]
  \centering
  \scriptsize
  \resizebox{\columnwidth}{!}{\begin{tabular}{@{}llrrrrrl@{}}
  \toprule
  Suite & Contrast & $\Delta$Loss & $\Delta$Acc & $\Delta$ECE & $\Delta$Brier & $\Delta$Tail & Harm counts \\
  \midrule
  Camelyon17 ERM & workflow: proxy - baseline & +0.055 [+0.046,+0.066] & +0.001 [-0.001,+0.003] & +0.0321 [+0.0305,+0.0334] & +0.0027 [+0.0013,+0.0038] & +5.50 [+5.00,+6.04] & loss/ECE/tail 10/10,10/10,10/10; Brier 8/10 \\
  Camelyon17 Finetune & workflow: proxy - baseline & +0.619 [+0.328,+0.916] & +0.025 [-0.021,+0.071] & +0.0214 [+0.0177,+0.0248] & +0.0006 [-0.0034,+0.0051] & +6.23 [+3.46,+9.06] & loss/ECE/tail 9/10,10/10,9/10; Brier 4/10 \\
  \midrule
  Camelyon17 ERM & same trajectory: proxy - val-loss & +0.020 [+0.011,+0.030] & +0.003 [+0.002,+0.005] & +0.0068 [+0.0047,+0.0090] & -0.0011 [-0.0023,+0.0001] & +2.01 [+1.47,+2.60] & loss/ECE/tail 10/10,10/10,10/10; Brier 3/10 \\
  Camelyon17 Finetune & same trajectory: proxy - val-loss & +0.010 [-0.135,+0.145] & +0.031 [+0.008,+0.055] & -0.0026 [-0.0081,+0.0029] & -0.0050 [-0.0118,+0.0008] & +3.42 [+0.75,+6.13] & loss/ECE/tail 2/10,1/10,4/10; Brier 1/10 \\
  \bottomrule
\end{tabular}
}
  \caption{
    Camelyon standard-metric triangulation and same-trajectory selector-only contrasts.
    Workflow rows compare the proxy-selected suppressive checkpoint with the baseline-selected undistorted checkpoint.
    Same-trajectory rows compare proxy-best with validation-loss-best on the same suppressive P95 trajectory and do not use fallback.
    $\Delta$Tail is teacher-difficulty tail CVaR, used here as a fixed-bank stress diagnostic.
    For Camelyon17 finetune rows, ECE/Brier are validation-side calibration diagnostics from the saved finetune selector-calibration artifact; the ERM rows use held-out/test calibration.
    Workflow-row intervals are regenerated in this standard-metric triangulation artifact; small endpoint differences from Table~\ref{tab:core-case-uncertainty} reflect separate bootstrap resampling of the same seed-matched deltas.
  }
  \label{tab:camelyon-selector-diagnostics}
\end{table}

Table~\ref{tab:camelyon-erm-selector-dose} extends the Camelyon17 ERM same-trajectory selector diagnostic beyond the P95 stress setting.
Holding the suppressive objective and trajectory fixed, proxy-best selection is worse than validation-loss-best selection at P95 and P97 on held-out loss, ECE, and the fixed-bank teacher-tail diagnostic.
At P99 the result attenuates and becomes mixed: standard loss and Brier no longer worsen, but ECE and teacher-tail still do.
This is why we treat P95/P97 as clean selector-only evidence and P99 as an attenuation check rather than as a contradiction.

\begin{table}[h]
  \centering
  \scriptsize
  \resizebox{\columnwidth}{!}{\begin{tabular}{@{}lrrrrll@{}}
  \toprule
  Regime & $\Delta$Loss & $\Delta$ECE & $\Delta$Brier & $\Delta$Tail & Harm counts & Note \\
  \midrule
  P95 & +0.020 [+0.011,+0.030] & +0.0068 [+0.0047,+0.0091] & -0.0011 [-0.0024,+0.0002] & +2.01 [+1.46,+2.61] & loss/ECE/tail 10/10,10/10,10/10; Brier 3/10 & clean \\
  P97 & +0.023 [+0.018,+0.028] & +0.0105 [+0.0088,+0.0122] & -0.0013 [-0.0025,-0.0002] & +2.39 [+2.05,+2.71] & loss/ECE/tail 10/10,10/10,10/10; Brier 2/10 & clean \\
  P99 & -0.003 [-0.006,-0.001] & +0.0050 [+0.0029,+0.0071] & -0.0023 [-0.0035,-0.0012] & +0.71 [+0.42,+1.02] & loss/ECE/tail 1/10,8/10,8/10; Brier 0/10 & attenuated/mixed \\
  \bottomrule
\end{tabular}
}
  \caption{
    Camelyon17 ERM same-trajectory selector-dose diagnostic.
    Deltas are proxy-best minus validation-loss-best on the same suppressive trajectory, so this isolates selector choice within a fixed objective/trajectory rather than a baseline-referenced workflow comparison.
    P95 and P97 show clean selector-only harm on held-out loss, ECE, and the fixed-bank teacher-tail diagnostic.
    P99 is attenuated and mixed: loss and Brier no longer worsen, but ECE and teacher-tail still degrade.
  }
  \label{tab:camelyon-erm-selector-dose}
\end{table}

Table~\ref{tab:camelyon-valacc-selector} checks validation-accuracy selection, a common scalarized checkpoint rule.
The comparison is validation-accuracy-best minus validation-loss-best on the same Camelyon17 ERM suppressive trajectory.
It is not a new workflow-generalization claim; it shows that the same practical selector warning is not limited to the soft-clipped proxy name.

\begin{table}[h]
  \centering
  \scriptsize
  \resizebox{\columnwidth}{!}{\begin{tabular}{@{}lrrrrll@{}}
  \toprule
  Regime & $\Delta$Loss & $\Delta$ECE & $\Delta$Brier & $\Delta$Tail & Harm counts & Note \\
  \midrule
  P95 & +0.023 [+0.017,+0.030] & +0.0072 [+0.0060,+0.0084] & -0.0020 [-0.0031,-0.0009] & +2.39 [+2.08,+2.72] & loss/ECE/tail 10/10,10/10,10/10; Brier 2/10 & clean loss/ECE/tail \\
  P97 & +0.027 [+0.023,+0.030] & +0.0117 [+0.0103,+0.0132] & -0.0027 [-0.0037,-0.0018] & +2.92 [+2.70,+3.19] & loss/ECE/tail 10/10,10/10,10/10; Brier 0/10 & clean loss/ECE/tail \\
  P99 & +0.026 [+0.019,+0.031] & +0.0221 [+0.0178,+0.0256] & +0.0002 [-0.0008,+0.0012] & +2.90 [+2.36,+3.33] & loss/ECE/tail 10/10,10/10,10/10; Brier 5/10 & clean loss/ECE/tail \\
  \bottomrule
\end{tabular}
}
  \caption{
    Camelyon17 ERM validation-accuracy selector sensitivity.
    Deltas are validation-accuracy-best minus validation-loss-best on the same suppressive trajectory.
    P95, P97, and P99 all worsen held-out loss, ECE, and the fixed-bank teacher-tail diagnostic in $10/10$ seeds.
    Brier is neutral or favorable to validation-accuracy selection, so the claim is restricted to loss, calibration by ECE, and the teacher-tail stress diagnostic.
  }
  \label{tab:camelyon-valacc-selector}
\end{table}

Table~\ref{tab:camelyon-loo-selector-standard} repeats a standard-metric selector-only check over leave-one-hospital-out Camelyon17 ERM folds.
Because the comparison is proxy-best minus validation-loss-best on the same suppressive trajectory, it supports selector-only robustness.
We do not use it to claim that proxy-minus-baseline workflow harm is uniform across hospitals.

\begin{table}[h]
  \centering
  \scriptsize
  \resizebox{\columnwidth}{!}{\begin{tabular}{@{}lrrrrl@{}}
  \toprule
  Regime & $\Delta$Loss & $\Delta$Acc & $\Delta$Brier & $\Delta$ECE & Harmful folds \\
  \midrule
  P95 & +0.020 & -0.0030 & +0.0028 & +0.0057 & loss/Brier 5/5,5/5; acc/ECE 5/5,5/5 \\
  P97 & +0.011 & -0.0014 & +0.0018 & +0.0040 & loss/Brier 5/5,5/5; acc/ECE 4/5,4/5 \\
  P99 & +0.008 & -0.0014 & +0.0019 & +0.0067 & loss/Brier 5/5,5/5; acc/ECE 5/5,5/5 \\
  \bottomrule
\end{tabular}
}
  \caption{
    Leave-one-hospital-out Camelyon17 ERM selector-only standard-metric check.
    Deltas are proxy-best minus validation-loss-best on the same suppressive trajectory, averaged over five held-out hospital folds.
    These rows use standard held-out metrics only and do not use the teacher-tail diagnostic.
    They are included as a selector-only robustness check, not as evidence that the baseline-referenced proxy-minus-baseline workflow hazard generalizes uniformly across hospitals.
  }
  \label{tab:camelyon-loo-selector-standard}
\end{table}

\subsection{Objective-Family Numerical Controls}
Table~\ref{tab:mechanism-table} gives the numerical values behind Figure~\ref{fig:mechanism}.

\begin{table}[h]
  \centering
  \small
  \resizebox{\columnwidth}{!}{\begin{tabular}{lcccc}
  \toprule
  Setting & $\Delta$Proxy$\downarrow$ & $\Delta$Tail CVaR$\downarrow$ & $\Delta$Held-out acc$\uparrow$ & $R_w$ \\
  \midrule
  Focal 1 & -0.055 & -2.435 & -0.0016 & 2.417 \\
  Focal 2 & -0.066 & -3.342 & -0.0018 & 4.100 \\
  Focal 4 & -0.070 & -4.113 & -0.0035 & 7.485 \\
  GCE q=0.7 & -0.024 & +5.843 & +0.0039 & 0.339 \\
  LS 0.02 & -0.025 & -1.046 & -0.0005 & 1.078 \\
  LS 0.05 & -0.029 & -1.768 & -0.0016 & 1.164 \\
  LS 0.10 & -0.033 & -2.444 & -0.0022 & 1.286 \\
  LS 0.20 & -0.040 & -3.271 & -0.0019 & 1.499 \\
  SoftClip P95 & -0.224 & +5.501 & +0.0010 & 0.200 \\
  SoftClip P97 & -0.098 & +5.155 & -0.0009 & 0.265 \\
  SoftClip P99 & +0.063 & +1.765 & -0.0009 & 0.605 \\
  \bottomrule
\end{tabular}
}
  \caption{
    Same-stack objective-family controls on Camelyon17 ERM.
    Soft clipping suppresses the tail ($R_w<1$), while focal loss and label smoothing upweight it ($R_w>1$).
    Focal 1/2/4 denote focal-loss focusing parameter $\gamma=1,2,4$.
    The $R_w$ values here are selected-checkpoint orientation values; Appendix Table~\ref{tab:rw-predictive-summary} and Figure~\ref{fig:rw-predictive-summary-app} separately report early-window $R_w$ for the warning-statistic analysis.
  }
  \label{tab:mechanism-table}
\end{table}

\subsection{Hard-Bootstrap Corroboration}
Table~\ref{tab:bootstrap-camelyon-support} gives the fuller second-family selector study on the same Camelyon17 ERM stack using hard bootstrapped cross-entropy with $\beta=0.60$, a standard suppressive objective from the noisy-label literature.
The core selector mismatch survives this extension.
Relative to val-loss-only selection inside the same bootstrap runs, proxy-only improves the bootstrap proxy by about $0.010$ but worsens held-out hospital-2 loss by $+0.0109$, teacher-difficulty tail CVaR by $+1.37$, and ECE by $+0.0110$ on average across ten seeds.
The selected checkpoints remain suppressive ($R_w \approx 0.567$, early $R_w \approx 0.605$).

The fixed $1.25\times$ guardrail is effectively slack on this family, while the within-run oracle-loss selector lands much earlier (mean epoch $6.4$ versus $15.2$ for proxy-only) and is lower on held-out loss, teacher-difficulty tail CVaR, and ECE.
Unlike the main P95 soft-clipping anchor, this bootstrap corroboration is not a stronger fixed-horizon pathology: relative to fixed epoch 30, the selected bootstrap checkpoint is actually better on proxy, held-out loss, teacher-difficulty tail CVaR, and ECE.
So we use it as supporting evidence against the ``only soft clipping'' objection rather than as a replacement for the main hazard anchors.

\begin{table}[h]
  \centering
  \small
  \resizebox{0.92\columnwidth}{!}{\begin{tabular}{lcccc}
  \toprule
  Selector & $\Delta$Held-out loss$\downarrow$ & $\Delta$Held-out acc$\uparrow$ & $\Delta$Tail CVaR$\downarrow$ & $\Delta$ECE$\downarrow$ \\
  \midrule
  Baseline & +0.000 & +0.000 & +0.000 & +0.000 \\
  Proxy-only & +0.016 & -0.003 & +1.892 & +0.008 \\
  Val-loss-only & +0.005 & -0.002 & +0.519 & -0.003 \\
  1.25x guardrail & +0.016 & -0.003 & +1.892 & +0.008 \\
  Oracle loss & -0.000 & +0.003 & +0.921 & +0.001 \\
  \bottomrule
\end{tabular}
}
  \caption{
    Hard-bootstrap selector extension on Camelyon17 ERM ($\beta=0.60$, ten seeds).
    Deltas are measured relative to the baseline-selected checkpoint.
    Proxy-only and val-loss-only are selectors on the same bootstrap runs; oracle loss is the best held-out-loss epoch within those same runs.
    The selector-level hazard is still the proxy-only versus safer-selector gap:
    proxy-only is slightly better on its own bootstrap proxy but materially worse on held-out loss, teacher-difficulty tail CVaR, and ECE.
    The fixed $1.25\times$ guardrail is included for completeness and is effectively slack on this family.
  }
  \label{tab:bootstrap-camelyon-support}
\end{table}

\subsection{Generalized Cross-Entropy Corroboration}
Table~\ref{tab:gce-q07-extension} gives a second standard noisy-label objective check on the Camelyon17 ERM stack.
Generalized cross-entropy with $q=0.7$ is suppressive on this stack, and the objective-aligned proxy selector moves to later checkpoints than validation-loss selection.
The result is supporting evidence, not a new optimizer comparison: it shows that the soft-clipping hazard also appears for a standard suppressive objective family.
We use $q=0.7$ as a single mid-range corroborating setting rather than a tuned GCE sweep.

\begin{table}[h]
  \centering
  \scriptsize
  \resizebox{\columnwidth}{!}{\begin{tabular}{llrrrrrr}
\toprule
Contrast & Epochs & $\Delta$GCE proxy & $\Delta$Acc & $\Delta$Loss & $\Delta$ECE & $\Delta$Tail & $R_w$ \\
\midrule
GCE proxy-best $-$ ERM baseline-selected & 14.4 / 3.8 & -0.0243 & +0.0039 & +0.0449 [+0.0391,+0.0514] & +0.0274 [+0.0261,+0.0287] & +5.84 [+5.54,+6.17] & 0.339 \\
GCE proxy-best $-$ GCE val-loss-best & 14.4 / 3.3 & -0.0146 & +0.0098 & +0.0181 [+0.0136,+0.0233] & +0.0129 [+0.0104,+0.0153] & +3.05 [+2.50,+3.60] & 0.339 \\
GCE fixed epoch 30 $-$ ERM fixed epoch 30 & 30 / 30 & -0.0173 & +0.0068 & +0.0648 [+0.0636,+0.0662] & +0.0308 [+0.0305,+0.0311] & +6.15 [+6.08,+6.21] & 0.308 \\
\bottomrule
\end{tabular}
}
  \caption{
    Generalized cross-entropy (GCE, $q=0.7$) extension on Camelyon17 ERM, ten seeds.
    Negative $\Delta$GCE proxy means the GCE objective-aligned validation score improves.
    The same-trajectory comparison is GCE proxy-best minus GCE validation-loss-best; all three contrasts worsen held-out loss, ECE, and teacher-tail in $10/10$ seeds.
    Epochs report mean selected epoch for the first and second policy in each contrast.
  }
  \label{tab:gce-q07-extension}
\end{table}

\subsection{Early Warning Value of $R_w$}
Table~\ref{tab:rw-predictive-summary} and Figure~\ref{fig:rw-predictive-summary-app} show that the same Camelyon17 ERM family-control stack already contains a practical early-warning signal.
Using the rule ``warn about harmful tail direction iff $R_w<1$,'' early $R_w$ computed well before final selection gives a strong same-stack warning about whether the eventual selected checkpoint will inflate or reduce the tail on the tested stack.

\begin{table}[h]
  \centering
  \small
  \resizebox{0.82\columnwidth}{!}{\begin{tabular}{lccc}
  \toprule
  Early window & Seed sign acc. & Regime sign acc. & Regime Spearman \\
  \midrule
  Epochs 1--3 & 0.97 & 1.00 & -0.964 \\
  Epochs 1--5 & 0.97 & 1.00 & -0.952 \\
  Epochs 1--10 & 0.97 & 1.00 & -0.952 \\
  \bottomrule
\end{tabular}
}
  \caption{
    Early $R_w$ as a same-stack warning for final selected-checkpoint tail direction on the Camelyon17 ERM objective-family stack.
    ``Regime'' aggregates the ten-seed mean for each family setting, and the sign rule warns about harmful tail direction when $R_w<1$.
    The warning is stable across short early windows.
  }
  \label{tab:rw-predictive-summary}
\end{table}

\begin{figure}[h]
  \centering
  \includegraphics[width=\columnwidth]{camelyon17_rw_predictive_summary_20260329}
  \caption{
    Early $R_w$ as a same-stack warning for final selected $\Delta$Tail on the Camelyon17 ERM objective-family stack.
    Each point is a ten-seed regime mean with $95\%$ mean confidence bars.
    The horizontal and vertical reference lines mark zero tail change and $R_w=1$ respectively.
  }
  \label{fig:rw-predictive-summary-app}
\end{figure}

\subsection{End-to-End Objective-Family Support}
Table~\ref{tab:objfam-finetune-support} summarizes the selected-checkpoint end-to-end Camelyon17 controls behind the mechanism discussion.
It is included to show that the same evaluation protocol yields a harmful suppressive family and comparison families with lower held-out loss and teacher-tail diagnostics.

\begin{table}[h]
  \centering
  \small
  \resizebox{\columnwidth}{!}{\begin{tabular}{lcccc}
  \toprule
  Family & $\Delta$Held-out loss$\downarrow$ & $\Delta$Held-out acc$\uparrow$ & $\Delta$Tail CVaR$\downarrow$ & $\Delta$ECE$\downarrow$ \\
  \midrule
  SoftClip P95 & +0.221 & +0.047 & +6.626 & +0.025 \\
  Label smoothing 0.10 & -0.276 & +0.029 & -2.647 & +0.014 \\
  Focal 2 & -0.280 & -0.026 & -2.526 & +0.024 \\
  \bottomrule
\end{tabular}
}
  \caption{
    Selected-checkpoint end-to-end objective-family support on a separate Camelyon17 finetune-control sweep.
    Deltas are measured relative to that sweep's own baseline finetune checkpoint.
    This table supports the method-ranking comparison in Table~\ref{tab:ranking-reversal} and is not numerically comparable to the main selector-anchor finetune sweep in Tables~\ref{tab:core-cases} and \ref{tab:selector-comparison}.
    Soft clipping is the harmful suppressive family; label smoothing is the cleanest counterexample; focal is mixed.
  }
  \label{tab:objfam-finetune-support}
\end{table}
